\chapter{Introducción: del Mpemba clásico al cuántico}
\label{ch:intro}

\section{El efecto Mpemba clásico (macroscópico)}

La observación de que, bajo condiciones idénticas de enfriamiento, una muestra de
agua inicialmente más caliente puede congelarse antes que una más fría fue
documentada por Mpemba y Osborne (1969). Durante décadas el fenómeno fue
controvertido por depender de detalles experimentales difíciles de controlar
(convección, evaporación, sobreenfriamiento, gases disueltos). El cambio de
paradigma llegó al reconocer que el ``efecto Mpemba'' no es una propiedad del
agua, sino una propiedad genérica de la \emph{relajación fuera del equilibrio}:
puede definirse de forma limpia en cualquier sistema cuya dinámica admita un
estado estacionario único y una noción de ``distancia al equilibrio''.

En el lenguaje markoviano clásico (Lu \& Raz, 2017), un sistema con
distribución de probabilidad $\bm{p}(t)$ relaja según $\dot{\bm{p}} = R\,\bm{p}$
hacia el equilibrio $\bm{\pi}$. La descomposición espectral de la matriz de tasas
$R$ controla las escalas de relajación, y el efecto Mpemba aparece cuando la
preparación caliente tiene un \emph{solapamiento anómalamente pequeño} con el
modo lento de $R$: relaja gobernada por modos rápidos y adelanta a la
preparación fría. Esta reformulación abstracta tuvo dos consecuencias: permitió
observar el efecto en plataformas controladas (fluidos granulares, coloides en
trampas ópticas, vidrios de espín) y abrió la puerta a su análogo cuántico.

\section{El efecto Mpemba cuántico}

En un sistema cuántico abierto markoviano la dinámica del operador densidad
$\rho$ está gobernada por la ecuación maestra de Gorini--Kossakowski--Sudarshan--Lindblad
(GKSL):
\begin{equation}
    \frac{\dd\rho}{\dd t} = \Lop[\rho]
    = -i[H,\rho] + \sum_{\mu}\left( L_\mu\rho L_\mu^\dagger
      - \tfrac{1}{2}\{L_\mu^\dagger L_\mu, \rho\}\right),
    \label{eq:gksl}
\end{equation}
con $\hbar=1$. El generador lineal $\Lop$ (el \emph{Liouvilliano}) es completamente
positivo y preserva la traza, y su espectro determina las escalas de relajación.
Suponiendo $\Lop$ diagonalizable, la dinámica admite la representación espectral
\begin{equation}
    \rho_t = \rss + \sum_{k\ge 2} e^{\lambda_k t}\,
             \underbrace{\Tr(\hat l_k^\dagger \rho_0)}_{a_k}\,\hat r_k,
    \label{eq:spectral}
\end{equation}
donde $\lambda_k$ son los autovalores de $\Lop$ (con $\Re\lambda_k\le 0$,
$\lambda_1=0$ el estado estacionario), y $\hat r_k$, $\hat l_k$ los autooperadores
derecho e izquierdo. El \textbf{efecto Mpemba cuántico} (QMpE) ocurre cuando una
preparación más alejada del equilibrio relaja antes porque su solapamiento $a_2$
con el \emph{modo lento} $\hat r_2$ es anómalamente pequeño (o nulo, en el caso
``fuerte'' $a_2=0$). El Liouvilliano es genéricamente \emph{no normal}
($\Lop\Lop^\dagger\neq\Lop^\dagger\Lop$): sus autooperadores no son ortogonales, y
es esa asimetría la que hace posible el efecto.

\section{Por qué exige cómputo de alto rendimiento}

El interés del QMpE no es solo conceptual: permite acelerar exponencialmente la
convergencia al estacionario mediante una preparación adecuada, con aplicaciones
en inicialización de qubits y termometría cuántica. La contrapartida es el coste:
para $N$ qubits, el espacio de operadores escala como $d^2 = 4^N$, de modo que el
Liouvilliano vive en dimensión $4^N\times 4^N$. La diagonalización densa cuesta
$\mathcal{O}(d^6)=\mathcal{O}(64^N)$ y es inviable más allá de $N\approx 6$. Esto
motiva una formulación numérica orientada a HPC, que el informe base descompone en
tres tareas computacionales (\S 7.3):

\begin{itemize}
  \item \textbf{T1 --- Espectro y modos lentos.} Extraer los pocos autovalores de
  $\Lop$ más cercanos a $0$ ($\lambda_2,\lambda_3$) y sus autooperadores, para
  diagnosticar el efecto. Problema de autovalores \emph{interior}.
  \item \textbf{T3a --- Trayectorias cuánticas.} Régimen de muchos cuerpos por
  saltos cuánticos: propagar vectores de estado $\ket{\psi}\in\mathbb{C}^d$ en vez
  del operador densidad ($\mathbb{C}^{d\times d}$). Reduce la memoria de $4^N$ a
  $2^N$.
  \item \textbf{T3b --- Redes tensoriales.} Representar el operador densidad como
  un MPS en el espacio doblado y evolucionarlo con TEBD disipativo, con coste
  polinómico en la dimensión de enlace $\chi$.
\end{itemize}

\section{Tesis del trabajo: por qué estas tres tareas}
\label{sec:tesis}

Conviene fijar desde el inicio la lógica que articula el trabajo, porque no es la
que podría suponerse. \textbf{El objetivo no es ``demostrar'' que el efecto Mpemba
existe} ---eso está establecido teórica y experimentalmente en la literatura
citada. El objetivo, siguiendo al informe base, es \emph{migrar la idea al régimen
cuántico construyendo las herramientas numéricas que permiten verificar y explotar
el QMpE en HPC}.

La clave está en una observación estructural: el QMpE es, en su núcleo, un
fenómeno de \emph{geometría espectral} de un generador no normal
(\cref{eq:gksl,eq:spectral}). Todo el fenómeno está codificado en dos objetos:
\begin{enumerate}
  \item la \textbf{estructura de relajación} ---los modos lentos $\lambda_2,\hat r_2,
  \hat l_2$ y los solapamientos $a_k=\Tr(\hat l_k^\dagger\rho_0)$--- cuyo valor
  \emph{predice} si una preparación puede exhibir el efecto ($a_2\!\approx\!0$ es
  la condición de Mpemba fuerte);
  \item la \textbf{dinámica de relajación} ---las curvas de distancia
  $\mathcal{D}(\rho_t\|\rss)$--- cuyo cruce \emph{es} la firma del efecto.
\end{enumerate}
Por tanto, ``verificar y explotar el QMpE'' se reduce a poder \emph{calcular} estos
objetos de forma correcta y escalable. Eso es exactamente lo que aportan las tres
tareas, y de ahí su elección:

\begin{itemize}
  \item \textbf{T1} calcula la estructura de relajación (modos lentos y $a_k$): es
  el \emph{diagnóstico} que predice el efecto. Pero requiere el Liouvilliano
  explícito, viable solo en pocos cuerpos.
  \item \textbf{T3a} y \textbf{T3b} calculan la dinámica de relajación en el
  régimen de \emph{muchos cuerpos} ---donde el fenómeno es físicamente más rico y
  el Liouvilliano explícito es imposible ($4^N$)--- por dos vías complementarias
  (trayectorias y redes tensoriales).
\end{itemize}

En consecuencia, \emph{cumplir} cada tarea significa demostrar dos cosas con datos:
(i) que el cálculo reproduce la física exacta (\textbf{verificación}, contra el
oráculo), y (ii) que escala en paralelo (\textbf{habilitación para explotar}, que
es el aporte de Sistemas Distribuidos). Como se mostrará, los datos lo confirman:
T1 extrae los modos lentos con error $\le10^{-5}$ y con aceleración algorítmica
\emph{faster-than-the-clock} ($8.8\times$); T3a reproduce la ecuación maestra con
la ley $M^{-1/2}$ y escala casi idealmente; y T3b converge al exacto
($\sim10^{-7}$) alcanzando $N=32$. La cadena queda así cerrada: las primitivas
que predicen y simulan el QMpE existen, son correctas y son escalables.

\section{Alcance y objetivos de este trabajo}

Este documento aborda las tareas \textbf{T1} y \textbf{T3} (a y b). Para cada una
se presenta: (i) el marco teórico, (ii) el método numérico y su discretización,
(iii) la implementación \emph{serial}, (iv) la justificación de la estrategia de
paralelización, y (v) la comparación y resultados (validación contra soluciones
exactas y escalado paralelo). La implementación es \emph{híbrida}: un núcleo de
referencia en Python (validado a precisión de máquina) y núcleos paralelos en
C++ con MPI + OpenMP (T1, T3a) y en Python con paralelismo de hilos (T3b). Todas
las cifras de validación y escalado de este documento provienen de ejecuciones
reales del código incluido, sobre la cadena de Ising transversa disipativa
\begin{equation}
    H = -J\sum_{i} Z_i Z_{i+1} - h\sum_i X_i, \qquad
    L_i^{\downarrow}=\sqrt{\gamma(\bar n+1)}\,\sigma_-^i,\quad
    L_i^{\uparrow}=\sqrt{\gamma\bar n}\,\sigma_+^i,
    \label{eq:ising}
\end{equation}
con $\bar n = (e^{2h/T}-1)^{-1}$ la ocupación de Bose del baño, que es el caso
paradigmático de muchos cuerpos del informe (\S 6.3).

\medskip
\noindent\textbf{Validación cruzada.} Los tres métodos se contrastan contra un
mismo \emph{oráculo}: la integración densa de \cref{eq:gksl} y la diagonalización
completa del Liouvilliano (módulo \texttt{common/qmpe.py}), que es exacta pero
solo viable para $N\lesssim 4$. Esto fija una base común de correctitud antes de
escalar.
